\chapter*{KATA PENGANTAR}
\addcontentsline{toc}{chapter}{KATA PENGANTAR}

Puji dan syukur penulis panjatkan ke hadirat Allah SWT atas rahmat, karunia, dan kemudahan yang diberikan sehingga penulis dapat menyelesaikan Laporan Tugas Akhir yang berjudul ``Pengembangan Aplikasi Pencarian Peluang Kolaborasi Mahasiswa: Modul Rekomendasi \emph{People-to-Project} Menggunakan \emph{Affinity Based Matching}''. Laporan ini disusun sebagai salah satu syarat untuk menyelesaikan program Sarjana di Program Studi Sistem dan Teknologi Informasi, Sekolah Teknik Elektro dan Informatika, Institut Teknologi Bandung.

Penulis menyadari bahwa penyelesaian Tugas Akhir ini tidak lepas dari bantuan, bimbingan, dan dukungan berbagai pihak. Oleh karena itu, pada kesempatan ini penulis ingin menyampaikan ucapan terima kasih yang sebesar-besarnya kepada pihak-pihak berikut.

\begin{enumerate}
    \item Allah SWT, atas segala rahmat, kekuatan, dan kelancaran yang diberikan sepanjang proses pengerjaan Tugas Akhir ini.

    \item Kedua orang tua penulis, Slamet Sutrisno dan Ulfah Hidayati, atas doa, kasih sayang, dukungan moral, dan dukungan materiil yang tiada henti selama penulis menempuh pendidikan hingga menyelesaikan Tugas Akhir ini.

    \item Kakak penulis, Alifah Salma Putri, atas doa, perhatian, dukungan, serta semangat yang diberikan selama proses perkuliahan dan penyusunan Tugas Akhir ini.

    \item Keluarga besar penulis yang senantiasa memberikan doa, dukungan, dan motivasi kepada penulis selama menempuh pendidikan.

    \item Bapak Dr. Ir. Arry Akhmad Arman, M.T. selaku dosen pembimbing yang telah meluangkan waktu, memberikan arahan, masukan, dan bimbingan yang sangat berharga selama proses penyusunan Tugas Akhir ini.

    \item Ir. I Gusti Bagus Baskara Nugraha, S.T., M.T., Ph.D. selaku Ketua Program Studi Sistem dan Teknologi Informasi, atas dukungan kelembagaan yang diberikan selama masa perkuliahan.

    \item Ahmad Fawwazi dan Muhammad Faishal Firdaus selaku rekan satu tim atas kerja sama dan koordinasi yang baik dalam menyatukan ketiga modul menjadi satu \emph{platform} yang terintegrasi.

    \item Teman-teman dekat penulis yang telah memberikan dukungan, semangat, dan bantuan selama proses pengerjaan Tugas Akhir ini.

    \item Seluruh dosen dan staf Program Studi Sistem dan Teknologi Informasi ITB yang telah memberikan ilmu dan pengalaman berharga selama masa perkuliahan.
\end{enumerate}

Penulis menyadari bahwa Tugas Akhir ini masih jauh dari sempurna. Oleh karena itu, penulis membuka diri terhadap kritik dan saran yang membangun untuk perbaikan di masa mendatang. Akhir kata, penulis berharap Tugas Akhir ini dapat memberikan manfaat bagi pembaca dan menjadi kontribusi yang berarti bagi pengembangan ilmu pengetahuan, khususnya di bidang sistem rekomendasi dan kolaborasi mahasiswa.

\begin{flushright}
    Bandung, Juli 2026\\

    Penulis \\
    \vspace{1cm}
    Muhammad Faishal Putra
\end{flushright}
